\documentclass{slate-resume}

\begin{document}

\slatemeta{John Doe Resume}{John Doe}{Resume}{robotics, software engineering}
\slateheader{John Doe}{1-555-000-0000\sep\href{mailto:john@example.com}{john@example.com}}{\href{https://www.example.com}{www.example.com}}

\section{Education}
\begin{entryblock}
\educationentry{Master of Science, Robotics}{Aug 2024 - May 2026}{Carnegie Mellon University, School of Computer Science}
\end{entryblock}
\entrygap
\begin{entryblock}
\educationentry{Bachelor of Science, Mechanical Engineering}{Aug 2020 - May 2024}{Georgia Institute of Technology, College of Engineering}
\end{entryblock}

\section{Relevant Experience}
\begin{entryblock}
\experienceentry{Robotics Software Engineer Intern}{Boston Dynamics}{Waltham, MA}{May 2025 - Aug 2025}
\begin{itemize}
  \item Implemented real-time obstacle avoidance algorithms for quadruped locomotion across uneven terrain, reducing collision events by 34\% during field trials in warehouse environments.
  \item Designed a ROS2-based sensor fusion pipeline integrating LiDAR, stereo vision, and IMU data to improve localization accuracy from 15cm to under 3cm drift per kilometer.
  \item Wrote comprehensive unit and integration tests for the perception stack, increasing code coverage from 52\% to 89\% and catching three critical race conditions before deployment.
  \item Presented weekly progress reports to a cross-functional team of 20, translating technical findings into actionable insights for the hardware and controls groups.
\end{itemize}
\end{entryblock}
\entrygap
\begin{entryblock}
\experienceentry{Undergraduate Research Assistant}{Georgia Tech LIDAR Lab}{Atlanta, GA}{Jan 2023 - May 2024}
\begin{itemize}
  \item Built a custom dataset of 12,000 annotated point clouds for outdoor scene segmentation, establishing a benchmark adopted by three partnering institutions.
  \item Trained transformer-based 3D object detection models on multi-GPU clusters, achieving a 7\% improvement in mean average precision over baseline architectures.
  \item Authored internal documentation and reproducible training scripts used by 15 lab members across two research groups, standardizing experiment tracking with Weights and Biases.
\end{itemize}
\end{entryblock}
\entrygap
\begin{entryblock}
\experienceentry{Autonomous Systems Intern}{Waymo}{Mountain View, CA}{May 2021 - Aug 2021}
\begin{itemize}
  \item Prototyped a lane-change prediction model using LSTM networks trained on 500 hours of highway driving data, improving prediction horizon accuracy by 18\% over the existing baseline.
  \item Automated nightly regression testing for the perception pipeline using Python and Bazel, cutting manual QA time by 3 hours per day across the team.
\end{itemize}
\end{entryblock}


\section{Notable Projects}
\begin{entryblock}
\projectentry{Atlas}{Autonomous drone navigation system with visual SLAM.}{C++, Python}
\begin{itemize}
  \item Engineered a monocular visual SLAM pipeline capable of mapping indoor environments at 30fps on embedded NVIDIA Jetson hardware with sub-5cm positional error.
  \item Integrated a lightweight path planner that dynamically reroutes around newly detected obstacles within 200ms, tested across 50 unique indoor floorplans.
\end{itemize}
\end{entryblock}
\entrygap
\begin{entryblock}
\projectentry{Forge}{Custom 3D printing slicer optimized for multi-material extrusion.}{Rust}
\begin{itemize}
  \item Developed a parallel slicing engine that processes complex multi-body STL files 3x faster than conventional open-source alternatives while maintaining sub-layer accuracy.
  \item Implemented automatic support structure generation using voxel-based overhang detection, reducing failed prints by 22\% across a test suite of 200 models.
\end{itemize}
\end{entryblock}



\section{Technical Skills}
\begin{entryblock}
\skillentry{Languages}{Python, C++, Rust, Go, TypeScript, MATLAB, SQL, Bash, \LaTeX}
\skillentry{Frameworks}{ROS2, PyTorch, TensorFlow, OpenCV, React, Docker, Kubernetes}
\skillentry{Hardware}{NVIDIA Jetson, Raspberry Pi, Arduino, STM32, ZED Stereo Camera, Velodyne LiDAR}
\end{entryblock}

\end{document}
